\documentclass[11pt]{article}
\usepackage[margin=1in]{geometry}
\usepackage{amsmath,amssymb}
\usepackage{graphicx}
\usepackage{xcolor}
\usepackage{hyperref}
\hypersetup{colorlinks=true, linkcolor=blue, citecolor=blue, urlcolor=blue}

\newcommand{\bP}{\mathbf{P}}
\newcommand{\bv}{\mathbf{v}}
\newcommand{\bu}{\mathbf{u}}
\newcommand{\bh}{\mathbf{h}}
\newcommand{\br}{\mathbf{r}}
\newcommand{\bp}{\mathbf{p}}
\newcommand{\bR}{\mathbf{R}}

\title{\textbf{Two Cells, One Closed Box, Merging}\\
\large Physics, Mechanics, and What Changed from Version 3}
\author{BTP — Phase-Field Modelling of Crawling Cells}
\date{\today}

\begin{document}
\maketitle

\begin{abstract}
\noindent
This note documents the physics and the control logic implemented in
\texttt{BTP\_version3\_two\_cell\_walk.ipynb}: two independently
crawling cells, each repeatedly choosing a direction, deforming to
move that way, relaxing back towards a circle, and choosing again,
confined to a closed box, that coalesce into one cell when their
paths cross. The underlying continuum model is exactly Version 3's
full phase-field / active-polar-fluid description with a spectral
Stokes solve; this note is organised around \emph{what changed} to
produce the new behaviour, and why each change was made the way it
was. \textbf{The notebook has not been executed as of this writing};
Section~\ref{sec:results} describes the designed, expected behaviour
and the diagnostics the code will produce, not measured output, and
says so explicitly wherever a number would otherwise look like a
result.
\end{abstract}

\tableofcontents

%======================================================================
\section{Scope and an upfront caveat}
%======================================================================

\texttt{BTP\_version3\_two\_cell\_walk.ipynb} was written and
statically checked (every cell compiles; a \texttt{pyflakes} pass
returns no warnings) but \textbf{deliberately not run}, per the
project owner's request to execute new code on their own machine.
Consequently this note cannot report measured speeds, merge times, or
rendered figures. What it \emph{can} do, and what it does below, is
explain precisely what the model computes, why each new mechanism was
added, how it differs from Version 3 and from the earlier Version 4
two-cell experiment, and what the diagnostics are designed to show
once the notebook is run. Every quantitative parameter below (box
size, wall strength, relaxation time, distance ranges) is a
\emph{chosen input}, not a fitted or verified one.

%======================================================================
\section{Recap: the Version 3 physics, unchanged}
\label{sec:v3recap}
%======================================================================

The two-cell notebook keeps the entire single-cell Version 3 model
intact; nothing in this section was modified. It is restated here
because every later section refers back to these equations by name.

\subsection{Fields}
A cell is a diffuse phase field $\phi(x,y,t)$, equal to $1$ inside the
cell and $0$ outside, with a smooth interface of width set by
$\varepsilon$. The actin cytoskeleton's local orientation and
polymerisation activity is a vector field $\bP(x,y,t)$, non-zero only
where $\phi>0$. The cytoplasm's velocity is $\bv(x,y,t)$, obtained at
every instant from a force balance (not integrated forward in time —
see Section~\ref{sec:nointertia}).

\subsection{Free energy}
Following Tjhung, Tiribocchi, Marenduzzo and Cates~\cite{Tjhung2015},
\begin{equation}
F[\phi,\bP] = \int d^2r\;\Big\{
\underbrace{a\,\phi^2(\phi-1)^2}_{\text{double well}}
+\underbrace{\tfrac{\varepsilon^2}{2}|\nabla\phi|^2}_{\text{interfacial tension}}
-\underbrace{\tfrac{\alpha}{2}(2\phi-1)|\bP|^2+\tfrac{\alpha}{4}|\bP|^4}_{\text{polar order inside the cell}}
+\underbrace{\tfrac{\kappa}{2}(\nabla\bP)^2}_{\text{Frank elasticity}}
+\underbrace{\beta\,\bP\cdot\nabla\phi}_{\text{anchoring}}
\Big\}.
\label{eq:F}
\end{equation}
The double well plus the interfacial term give the cell a boundary
with a cost; the $\alpha$ terms drive $|\bP|^2\to2\phi-1$ strictly
inside the cell; the Frank term penalises sharp turns in $\bP$; and
the anchoring term, with $\beta>0$, makes $\bP$ prefer to point
outward, normal to the local boundary — in every version of this
project, $\beta$ has been the dominant shape parameter.

Differentiating gives the chemical potential driving $\phi$ and the
molecular field driving $\bP$:
\begin{align}
\mu_\phi &= \frac{\delta F}{\delta\phi}
= 2a\,\phi(\phi-1)(2\phi-1) - \varepsilon^2\nabla^2\phi - \alpha|\bP|^2 - \beta\,\nabla\!\cdot\!\bP,
\label{eq:mu}\\
\bh &= \frac{\delta F}{\delta\bP}
= \big[-\alpha(2\phi-1)+\alpha|\bP|^2\big]\bP - \kappa\nabla^2\bP + \beta\nabla\phi.
\label{eq:h}
\end{align}

\subsection{Equations of motion}
With $\bu=\bv+w_0\bP$ the net velocity transporting cell material
($w_0$ the actin treadmilling rate),
\begin{align}
\partial_t\phi + \nabla\!\cdot\!(\phi\,\bu) &= -M\big[\mu_\phi-\Lambda\,g(\phi)\big],
\qquad g(\phi)=\phi(1-\phi),
\label{eq:phidot}\\
\partial_t\bP + (\bu\cdot\nabla)\bP &= -\boldsymbol\Omega\cdot\bP + \chi\,\mathbf D\cdot\bP - \bh/\Gamma,
\label{eq:Pdot}
\end{align}
where $\mathbf D$ and $\boldsymbol\Omega$ are the symmetric
(strain-rate) and antisymmetric (vorticity) parts of
$\partial_\alpha v_\beta$, $\Gamma$ a rotational viscosity, and
$\chi$ the flow-alignment parameter. $\Lambda$ is a spatially uniform
Lagrange multiplier, $\Lambda=\int\mu_\phi\,dA/\int g\,dA$, which makes
Eq.~\eqref{eq:phidot} conserve $\int\phi\,dA$ \emph{exactly}, to
machine precision, by construction — this is the
Rubinstein--Sternberg conservative form of Allen--Cahn
dynamics~\cite{Rubinstein1992}, not an approximate correction.

\subsection{Force balance}
At the cell's scale $Re\sim10^{-5}$, so inertia is dropped and the
force balance is the Stokes limit,
\begin{equation}
0 = -\nabla p + \eta\nabla^2\bv - \xi\bv + \mathbf f,
\qquad \nabla\!\cdot\!\bv = 0,
\label{eq:stokes}
\end{equation}
with the substrate contribution $-\xi\bv$ standing in for the wall of
the original 3D geometry (the continuum limit of a carpet of
stochastic adhesive linkers~\cite{LeoniSens2017}), and the body force
\begin{equation}
f_\alpha = \partial_\beta\!\big(\sigma^{\rm act}_{\alpha\beta}+\sigma^{\rm el}_{\alpha\beta}\big) - \phi\,\partial_\alpha\mu_\phi,
\qquad
\sigma^{\rm act}_{\alpha\beta} = -\zeta\,\phi\,P_\alpha P_\beta,
\label{eq:force}
\end{equation}
$\zeta<0$ contractile, $\sigma^{\rm el}$ the elastic (Frank $+$
co-rotational) stress, and the last term of Eq.~\eqref{eq:force} the
capillary force in potential form.

\subsection{The spectral solve}
Because the domain is periodic, Eq.~\eqref{eq:stokes} is solved
exactly in Fourier space: a Leray (transverse) projection eliminates
the pressure, and dividing by $\xi+\eta k^2$ gives $\hat{\bv}(\mathbf
k)$ directly — one pair of FFTs per step, no iteration, no time-step
restriction from the viscous term. \textbf{This spectral method is
the one piece of Version 3's numerics that only works because the
domain is periodic}, a fact Section~\ref{sec:wall} returns to.

%======================================================================
\section{What changed: from one cell to two, in a closed box}
\label{sec:changes}
%======================================================================

Table~\ref{tab:changes} gives the short version; the rest of this
section explains each row.

\begin{table}[h]
\centering
\small
\begin{tabular}{@{}p{3.6cm}p{5.6cm}p{5.6cm}@{}}
\hline
& \textbf{Version 3 (single cell)} & \textbf{Two-cell walk notebook} \\
\hline
Grid & $N{=}128$, $L{=}80$ & $N{=}192$, $L{=}120$ (room for two cells and a walk) \\
Domain & periodic, infinite in effect & periodic numerically, but confined by a new soft wall term \\
Phase/polarization fields & one $\phi$, one $\bP$, one cell & one $\phi$, one $\bP$, \emph{two disjoint blobs} at $t{=}0$ \\
Activity ($w_0,\zeta,\beta$) & uniform scalars, on for the whole run & $w_0$ becomes a \emph{field} (array); each blob's region is independently on/off \\
Direction of $\bP$ & set once at $t{=}0$, implicit in the initial condition & chosen explicitly, re-chosen on a schedule, per blob \\
What stops the cell at the edge & nothing (recentring trick keeps one cell near the middle) & the new wall potential (no recentring; absolute position matters) \\
Two cells touching & not modelled & coalesce into one connected $\phi$ region, detected via connected-component labelling \\
Post-collision direction & n/a & biased towards the two cells' combined pre-collision momentum \\
\hline
\end{tabular}
\caption{What changed, at a glance.}
\label{tab:changes}
\end{table}

\subsection{Grid and domain}
The grid grew from $128\times128$ ($L=80$) to $192\times192$
($L=120$) at the same resolution $dx=0.625$. Two cells need room to
start apart, walk, and potentially collide without ever nearing the
box edge closely enough for the soft wall (Section~\ref{sec:wall}) to
distort their shape before a genuine wall-avoidance decision should
fire.

\subsection{The soft confining wall}
\label{sec:wall}
Version 3's Stokes solve is spectral, which requires periodicity;
there is no way to impose a literal no-slip or free-slip wall inside
that method without replacing it with a different linear solve
(finite differences or finite volumes, discretising
Eq.~\eqref{eq:stokes} with wall boundary conditions and solving the
resulting sparse system — a substantial rewrite, not attempted here).
Instead, exactly \emph{one} new term is added to the chemical
potential of Eq.~\eqref{eq:mu}:
\begin{equation}
\mu_\phi \;\longrightarrow\; \mu_\phi + k_{\rm wall}\,W(x,y)\,\phi,
\qquad
W(x,y) = \left[\frac{\max\!\big(0,\,|x|-x_0\big)^2+\max\!\big(0,\,|y|-y_0\big)^2}{(L/2-x_0)^2}\right]^{1/2\cdot{\rm sharp}},
\label{eq:wall}
\end{equation}
precomputed once as a static array $W$: zero everywhere inside a
square region starting at $x_0=y_0=0.8\times(L/2)$ from the centre,
rising smoothly (quartic by default, set by \texttt{wall\_sharpness})
to $1$ at the literal domain edge. This is structurally the
\emph{same kind of term} as the pairwise steric-repulsion used in the
Version 4 two-cell experiment ($+2k_{\rm rep}\phi_i\phi_j^2$ in each
cell's $\mu_i$) — a quartic penalty that only activates where $\phi$
overlaps a forbidden region — except the forbidden region here is a
fixed wall field $W$ instead of another cell's $\phi_j$. It plugs into
the unmodified relaxation term of Eq.~\eqref{eq:phidot} with no other
code change.

\paragraph{What this is not.} It is a soft potential, not a boundary
condition. The fluid velocity field $\bv$ is still computed on a
periodic domain; what changes is that $\phi$ is discouraged, by an
energy penalty, from occupying the near-edge zone, which behaviourally
keeps the cell inside a confined region without literally terminating
the computational domain there. The strength $k_{\rm wall}=40$ and the
margin $0.8\times(L/2)$ are \emph{reasoned starting values}: they were
chosen by the same reasoning used for the Version 4 repulsion
coefficient (large enough, relative to the double-well coefficient
$a$, to resist the active crawling force pushing against it) but were
not empirically tuned by running this notebook. A runtime check
(\texttt{touches\_boundary}) prints a warning if $\phi>0.5$ ever
reaches within six grid cells of the literal array edge, which would
indicate $k_{\rm wall}$ needs to be raised.

\subsection{Two blobs, one field}
\label{sec:oneField}
The Version 4 two-cell experiment (a separate, earlier notebook) used
\emph{two separate} fields, $\phi_1,\bP_1$ and $\phi_2,\bP_2$, coupled
by a steric term specifically so the cells would stay distinguishable
and never overlap. This notebook does the opposite on purpose: it
uses \emph{one} $\phi$ and \emph{one} $\bP$, initialised with two
disjoint tanh profiles,
\begin{equation}
\phi(\br,0) = \max\Big(\tfrac12\big[1-\tanh(\tfrac{|\br-\br_1|-R_0}{\sqrt2\,\varepsilon})\big],\;
\tfrac12\big[1-\tanh(\tfrac{|\br-\br_2|-R_0}{\sqrt2\,\varepsilon})\big]\Big),
\end{equation}
and $\bP$ set independently within each blob's own disc to point
along that blob's chosen heading. While the two regions of $\phi>0$
are disjoint, nothing distinguishes them as separate objects to the
underlying PDE — Eqs.~\eqref{eq:phidot}--\eqref{eq:Pdot} are evaluated
over the whole grid every step, exactly as they are for one cell.
When the two regions' boundaries touch, the double-well and
interfacial-tension terms of Eq.~\eqref{eq:F} do what they always do:
minimise interfacial length, which for two adjacent blobs of the
\emph{same} field means rounding the pair into one connected shape.
\textbf{No merging algorithm was written}; coalescence is a direct
consequence of using one shared field, which is the most literal
reading of ``mix into one another'' available in a diffuse-interface
model — once the interface between them disappears, there is no
record left of which material came from which original cell.

\subsection{Activity as a field, not a scalar}
Before the two blobs merge, one of them may be actively crawling
(treadmilling, anchoring, active stress all switched on) while the
other is relaxing (all switched off), simultaneously, in the same
field. This requires only one structural change: $w_0$, previously a
single scalar multiplying $\bP$ in $\bu=\bv+w_0\bP$, becomes a
spatial array $w_0(x,y)$, equal to the active value inside whichever
blob is currently in its ``move'' phase and zero elsewhere (including
inside a blob that is relaxing). $\zeta$ and $\beta$ needed no such
change: they already appear multiplying other fields elementwise in
Eqs.~\eqref{eq:mu}--\eqref{eq:force}
(\texttt{chemical\_potential\_phi}, \texttt{molecular\_field\_P},
\texttt{body\_force}), so passing them in as arrays instead of scalars
works with the existing code unmodified — only the
function signature and the one line computing $\bu$ needed to accept
a field for $w_0$.

\subsection{The decision layer}
\label{sec:decision}
Each disjoint blob (and, after merging, the single resulting cell)
carries a small state: a current heading $\theta$, a phase (\emph{move}
or \emph{relax}), a target travel distance for the current move leg,
and the position where that leg began. Every \texttt{sample\_every}
steps ($=1$ time unit at $dt=0.01$), three conditions are checked, in
priority order:
\begin{enumerate}
\item \textbf{Near a wall.} If the blob's centroid is within
\texttt{wall\_redirect\_margin} of the box edge, pick a new heading
biased back toward the box centre (a von Mises draw centred on
$\arctan2(-y,-x)$, concentration \texttt{wall\_kappa}) and start a new
move leg immediately.
\item \textbf{Move leg finished.} If the blob is in \emph{move} phase
and has covered its target distance, switch to \emph{relax}: $w_0$,
$\zeta$, $\beta$ all set to zero within its mask. With no energy
sustaining the fan shape, Eq.~\eqref{eq:phidot}'s relaxation term
— the same interfacial relaxation validated in Version 3's
passive-limit test, an ellipse relaxing to a circle — pulls the blob
back towards a circle on its own; nothing new was added for this.
\item \textbf{Relax finished.} If the blob has been relaxing for at
least \texttt{relax\_time}, pick a fresh, uniformly random heading
and resume \emph{move}.
\end{enumerate}
Identifying which blob is which between checks uses
\texttt{scipy.ndimage.label} on $\phi>0.5$ to find connected
components, matched to the previous check's two tracked blobs by
whichever pairing minimises total centroid displacement — this keeps
a blob's identity (and hence its in-progress target distance and
heading) stable across checks even though the labelling algorithm
itself has no memory of which numbered label is which cell.

\subsection{The merge event}
\label{sec:merge}
The moment \texttt{ndimage.label} reports \textbf{one} connected
component where it previously reported two, the merge branch fires
exactly once. Using the \emph{last} check before contact (where the
two blobs were still separately resolvable), it computes each blob's
mass-like weight $m_i=\int_{\Omega_i}\phi\,dA$, its instantaneous
crawling velocity $\bv_i$ (centroid displacement over the preceding
check interval, divided by the elapsed time), and from these:
\begin{align}
\bp_i &= m_i\,\bv_i, &
\bp_{\rm tot} &= \bp_1+\bp_2,\\
\bR_{\rm cm} &= \frac{m_1\br_1+m_2\br_2}{m_1+m_2}, &
L_i &= m_i\,(\br_i-\bR_{\rm cm})\times\bv_i,
\qquad L_{\rm tot}=L_1+L_2.
\label{eq:momentum}
\end{align}
$L_i$ is the $z$-component of the orbital angular momentum each
blob's crawling motion carries about the eventual contact point —
exactly the quantity this project's adviser asked to be taken into
account. Section~\ref{sec:nointertia} explains why it is used as a
\emph{bias}, via a von Mises distribution centred on
$\arctan2(p_{{\rm tot},y},p_{{\rm tot},x})$ with concentration
\texttt{merge\_kappa}, for the merged cell's first new heading —
rather than enforced as a literal conserved quantity. $\bp_{\rm tot}$
and $L_{\rm tot}$ are both printed when the merge fires, as a
diagnostic record of what the pre-collision state implied, whether or
not the model's own stochastic choice ends up close to that
direction.

\subsection{Why not literal angular-momentum conservation}
\label{sec:nointertia}
This is worth stating plainly rather than quietly working around.
Stokes flow (Eq.~\eqref{eq:stokes}) is the zero-Reynolds-number
limit: inertia has been dropped from the equations entirely, which is
also why this project has used $Re\sim10^{-5}$ to justify the force
balance since Version 3 was first built. In that limit, momentum is
\emph{not} a quantity carried forward in time at all — net force and
net torque on the fluid are balanced \emph{instantaneously}, every
step, by viscous dissipation and the substrate friction $-\xi\bv$,
which is itself an external drag that can (and does) absorb momentum
from outside the fluid/cell system. There is no dynamical bookkeeping
in Eq.~\eqref{eq:stokes} analogous to $\tfrac{d}{dt}(m\bv)=0$ for an
isolated collision, because there is no inertial term to begin with.
This is the standard, if counter-intuitive, fact about swimming at low
Reynolds number: a cell that stops generating active stress stops
moving \emph{immediately}, it does not coast~\cite{Purcell1977}.

The notebook therefore does not — and physically should not — enforce
$\bp_{\rm tot}$ or $L_{\rm tot}$ as a literal post-collision
constraint on the merged cell's velocity or rotation; nothing in
Eq.~\eqref{eq:stokes} would support such a constraint without breaking
the force-balance structure the whole model rests on. What it does
instead, as described in Section~\ref{sec:merge}, is treat $\bp_{\rm
tot}$ as \emph{information}, biasing a stochastic decision the model
was always going to make (a new random heading) rather than
determining it outright. This keeps faith with the physics while
still giving the collision's approach geometry a say in what happens
next, which is the practical substance of the adviser's suggestion.

%======================================================================
\section{Mechanics of the control loop}
\label{sec:mechanics}
%======================================================================

Algorithmically, one iteration of the outer loop (one
\texttt{sample\_every} interval, $\approx1$ time unit):
\begin{enumerate}
\item Label $\phi>0.5$ into connected components.
\item \emph{If not yet merged and two components are found:} identify
and match them to the two tracked states (Section~\ref{sec:decision});
for each, evaluate the three triggers in order and, if one fires,
update that blob's state and repoint $\bP$ within its mask
(\texttt{set\_direction}); build the elementwise activity fields
$\zeta_{\rm run},\beta_{\rm run},w_0$ from both blobs' current phases.
\item \emph{If not yet merged and one component is found:} run the
merge computation of Eq.~\eqref{eq:momentum}, collapse to a single
tracked state with a momentum-biased new heading, repoint $\bP$ over
the whole merged region, and report $\bp_{\rm tot},L_{\rm tot}$.
\item \emph{If already merged:} evaluate the same three triggers for
the single state.
\item Advance \texttt{sample\_every} individual timesteps of
Eqs.~\eqref{eq:phidot}--\eqref{eq:Pdot} and \eqref{eq:stokes}
(\texttt{advance\_one\_step}), unchanged from Version 3 except that
$w_0$ is now an array; store a video frame every
\texttt{frame\_every} steps, and check for divergence.
\end{enumerate}
The loop ends when the elapsed simulated time reaches
\texttt{t\_final}.

%======================================================================
\section{Diagnostics and expected behaviour}
\label{sec:results}
%======================================================================

\textbf{No numbers in this section were measured — the notebook has
not been run.} This section describes what each diagnostic is
designed to show, and the qualitative behaviour the design is
intended to produce, so that a reader can judge the first real run
against it.

\begin{description}
\item[Path plot] Each tracked entity's centroid is logged every check
interval, tagged by identity (\texttt{0}, \texttt{1}, or
\texttt{merged}) and phase (\emph{move} or \emph{relax}). The intended
picture: two separate, roughly piecewise-linear trajectories (straight
during \emph{move}, nearly stationary during \emph{relax}) that
converge, mark a merge point, and continue as one trajectory
thereafter — visually distinguishing ``two independent walks'' from
``one combined walk'' without needing to inspect the field directly.
\item[Merge report] $\bp_{\rm tot}$, $L_{\rm tot}$, and the resulting
new heading, printed once. If the two cells approach nearly head-on
with similar masses, $\bp_{\rm tot}$ should be small and the new
heading close to uniformly random (the von Mises bias has little to
act on); if they approach at an angle with different masses or speeds,
$\bp_{\rm tot}$ should point roughly along the net approach direction
and bias the merged cell's first leg that way.
\item[Video] \texttt{two\_cell\_walk.mp4}, titled per frame with each
active entity's phase and (while moving) heading in degrees, intended
to make the whole decide/move/relax/redirect cycle — and the moment of
coalescence — legible without reading the printed log alongside it.
\item[Boundary warning] \texttt{touches\_boundary} should \emph{not}
fire during a normal run; if it does, $k_{\rm wall}$ or the wall
margin need to be increased (Section~\ref{sec:wall}).
\end{description}

\paragraph{Expected qualitative sequence.} Two blobs, placed
$\texttt{init\_separation}=50$ apart along $x$, each independently
cycle through move/relax legs of random length drawn from
$(15,35)$ and random heading, inside a box of half-width $60$ with
the wall zone beginning at $48$. Given random, independent headings,
whether and when the two paths cross is itself random — for a
\texttt{t\_final}$=400$ run there is no guarantee a merge event occurs
at all; the run should be repeated with a different \texttt{seed} or
a larger \texttt{t\_final}/smaller \texttt{init\_separation} if it
does not. The notebook prints an explicit message distinguishing a
run with a merge from one without, rather than assuming one will
happen.

%======================================================================
\section{Open questions and known limitations}
\label{sec:limitations}
%======================================================================

\begin{itemize}
\item \textbf{The wall is a potential, not a boundary condition}
(Section~\ref{sec:wall}); $k_{\rm wall}$ and the margin are untuned.
\item \textbf{Relaxation uses a fixed duration}, not the
circularity-based criterion of the single-cell run-and-tumble
notebook — simpler, at the cost of not actually checking the blob has
rounded off before resuming activity.
\item \textbf{The merge biases heading only.} A natural extension,
not implemented, is seeding a small swirl in $\bP$ proportional to the
sign of $L_{\rm tot}$, so the merged cell visibly starts to curve
rather than only picking a biased straight heading.
\item \textbf{Connected-component labelling can be confused by small
numerical artefacts.} The code defends against spurious extra
components by always tracking only the two largest, but has not been
stress-tested against this failure mode by execution.
\item \textbf{Generalising past two cells} needs the \texttt{states}
list and the nearest-centroid matching logic generalised from $2$ to
$N$ blobs; not attempted here.
\item \textbf{The honest long-term fix for Section~\ref{sec:wall}} is
a genuine non-periodic Stokes solve (a finite-difference
Poisson/Helmholtz solve with wall boundary conditions, replacing the
two-FFT spectral method with a sparse linear solve) rather than a soft
potential on a periodic domain.
\end{itemize}

%======================================================================
\begin{thebibliography}{9}
\bibitem{Tjhung2015}
E.~Tjhung, A.~Tiribocchi, D.~Marenduzzo, M.~E.~Cates,
\emph{A minimal physical model captures the shapes of crawling cells},
Nat. Commun. \textbf{6}, 5420 (2015).

\bibitem{LeoniSens2017}
M.~Leoni, P.~Sens,
\emph{Model of cell crawling controlled by mechanosensitive adhesion},
Phys. Rev. Lett. \textbf{118}, 228101 (2017).

\bibitem{Rubinstein1992}
J.~Rubinstein, P.~Sternberg,
\emph{Nonlocal reaction-diffusion equations and nucleation},
IMA J. Appl. Math. \textbf{48}, 249 (1992).

\bibitem{Purcell1977}
E.~M.~Purcell,
\emph{Life at low Reynolds number},
Am. J. Phys. \textbf{45}, 3 (1977).
\end{thebibliography}

\end{document}
